\documentclass[10pt]{article}
\usepackage[a4paper,margin=0.3in]{geometry}
\usepackage{etoolbox}
\AtBeginEnvironment{align}{\small}
\usepackage{amsmath,amssymb,amsfonts,mathtools}
\usepackage[dvipsnames]{xcolor}
\usepackage{array,booktabs,float,longtable}
\usepackage[colorlinks=true,linkcolor=blue]{hyperref}
\allowdisplaybreaks
\numberwithin{equation}{section}
\setlength{\jot}{3pt}
\setlength{\emergencystretch}{3em}
\definecolor{JIM}{RGB}{0,130,60}
\definecolor{CC}{RGB}{140,0,170}
\definecolor{BFK}{RGB}{215,95,0}
\newcommand{\new}[1]{{\color{blue}#1}}
\newcommand{\old}[1]{{\color{red}#1}}
\newcommand{\nc}{\color{blue}}
\newcommand{\rc}{\color{red}}
\newcommand{\NNv}{\mathcal N\log\frac{\vee}{k^+}}
\def\mf#1{\mathbb{#1}}

\title{Where the $\log\frac{\vee}{k^+}$ of Group III comes from:\\ step by step, with your Row VI}
\date{}
\begin{document}
\maketitle
\noindent Rows: your \texttt{Other\_terms.tex}. In your \texttt{BFKL.tex}, Group III Rows III, IV, V are ``NA''. The rows meant here are the rows of \texttt{Other\_terms.tex} whose line 2 has $\log(\frac{\vee}{k^+}-1)$: Group III Row III (both parts), Row V and Row VI (Row IV has only four $\rho$'s), and Group II Row IV. Notation as in your files.

\section{The answer in short}
Your Row VI is correct up to and including the plus prescription. After it, the only logarithm you can \emph{see} is $\big[\log\frac{k^+}{\Lambda}+\log(\frac{\vee}{k^+}-1)\big]$ in line 2. But line 1 is not independent of $\vee$: its upper limit is $\vee-k^+$. There are two different $\log\frac{\vee}{k^+}$ in the row:
\begin{enumerate}
\item \textbf{The visible one}, $\log(\frac{\vee}{k^+}-1)=\log\frac{\vee}{k^+}+\log(1-\frac{k^+}{\vee})$ in line 2. It is \emph{not} a BFKL term, and you are right to leave it out of \texttt{BFKL.tex}. It is cancelled exactly by the subtraction $-f(0)$ of $\left[\frac{1}{p^+}\right]_+$ in line 1, integrated up to $p^+=\vee-k^+$ (Step 5).
\item \textbf{A hidden one}, in line 1, in the term $-\frac{1}{k^+}\,\delta_{im}\,\delta_{jk'}$ of the bracket. As written it carries no logarithm, but its $p^+$ integral grows like $\log\frac{\vee}{k^+}$ for $\vee\to\infty$. This shows up only once you look at what the $y'$ integral does at large $p^+$ (Steps 3, 4). This one \emph{is} the BFKL term; it is called (B) (Step 6).
\end{enumerate}
So, row by row: Row VI $=\log\frac{k^+}{\Lambda}\times$(JIMWLK part) $+$ (B) $+$ a rest that is finite for $\vee\to\infty$ (Step 7). Without (B), the BFKL list misses the terms in which the measured gluon is emitted by the gluon at $z$: 9 of the 67 colour/kernel structures of the two-$\rho$ sum then differ from BFKL$\otimes$LO (Sec.~\ref{sec:checks}). Without $-$(B), the ``rest'' still grows like $\log\frac{\vee}{k^+}$.

\section{Step 1: your Row VI}
From \texttt{Other\_terms.tex} (your last equation for Row VI):
\begin{align}
    \frac{d^{3}N_{\rm NLO}}{d^{3}k} & = \frac{1}{(2\pi)^{3}}\, \frac{ig^{4}\,f^{c'd'a}}{4\pi^{4}}\, \frac{1}{k^{+}}\,\int_{0}^{\vee-k^{+}}\, \frac{dp^{+}}{2\pi}\, \int_{w,y',z,y,x,x'}\, e^{-ik\cdot (y'-w)}\, e^{-ik\cdot \frac{p^{+}}{k^{+}}\,(y'-z)}\,\frac{(y'-z)^{m}\,(x'-y')^{k'}\,(y-w)^{i}\,(x-z)^{j}}{(y'-z)^{2}\,(x'-y')^{2}\,(y-w)^{2}\,(x-z)^{2}}\notag \\ & \left[\frac{\delta_{k'm}\,\delta_{ij}}{(p^{+}+k^{+})}-\left[\frac{1}{p^{+}}\right]_{+}\,\delta_{jm}\,\delta_{ik'}-\frac{1}{k^{+}}\, \delta_{im}\, \delta_{jk'}\right]\, \left[U^{e'c'}(y')\,U^{d'b}(z)\,\rho^{e'}(x')-U^{bd'}(z)\,U^{c'e'}(x')\,\rho^{e'}(x')\right]\times\notag \\ &  \left[U^{bc}(z)\,U^{ad}(y)\,U^{ce}(x)\,\rho^{d}(y)\,\rho^{e}(x)-U^{ad}(y)\,\rho^{d}(y)\,\rho^{b}(x)\right]\notag \\ &  -\,\frac{1}{(2\pi)^{3}}\, \frac{ig^{4}\,f^{c'd'a}}{8\pi^{5}}\, \frac{1}{k^{+}}\,{\color{blue}\log{\left(\frac{\vee}{k^{+}}-1\right)}}\, \int_{w,y',z,y,x,x'}\, e^{-ik\cdot (y'-w)}\,\frac{(y'-z)\cdot (x-z)\,(x'-y')\cdot (y-w)}{(y'-z)^{2}\,(x'-y')^{2}\,(y-w)^{2}\,(x-z)^{2}}\notag \\ & \left[U^{e'c'}(y')\,U^{d'b}(z)\,\rho^{e'}(x')-U^{bd'}(z)\,U^{c'e'}(x')\,\rho^{e'}(x')\right]\times  \left[U^{bc}(z)\,U^{ad}(y)\,U^{ce}(x)\,\rho^{d}(y)\,\rho^{e}(x)-U^{ad}(y)\,\rho^{d}(y)\,\rho^{b}(x)\right] \label{eq:userG3VI}
\end{align}
Call the colour factor
\begin{equation}
\mathcal C(y')=\Big[U^{e'c'}(y')\,U^{d'b}(z)\,\rho^{e'}(x')-U^{bd'}(z)\,U^{c'e'}(x')\,\rho^{e'}(x')\Big]\times\Big[U^{bc}(z)\,U^{ad}(y)\,U^{ce}(x)\,\rho^{d}(y)\,\rho^{e}(x)-U^{ad}(y)\,\rho^{d}(y)\,\rho^{b}(x)\Big] .
\label{eq:C}
\end{equation}
Only the first bracket depends on $y'$, through $U^{e'c'}(y')$; $\mathcal C(z)$ is $\mathcal C(y')$ with $U^{e'c'}(y')\to U^{e'c'}(z)$.

\section{Step 2: line 1 is three terms}
Contract the Kronecker deltas of the three terms of the bracket of line 1 and name the transverse integrals:
\begin{align}
F_a(p^+)&=\int_{w,y',z,y,x,x'}\,e^{-ik\cdot (y'-w)}\, e^{-ik\cdot \frac{p^{+}}{k^{+}}\,(y'-z)}\,\frac{(y'-z)\cdot (x'-y')\,(y-w)\cdot (x-z)}{(y'-z)^{2}\,(x'-y')^{2}\,(y-w)^{2}\,(x-z)^{2}}\,\mathcal C(y') &&(\delta_{k'm}\delta_{ij}),\notag\\
F_b(p^+)&=\int_{w,y',z,y,x,x'}\,e^{-ik\cdot (y'-w)}\, e^{-ik\cdot \frac{p^{+}}{k^{+}}\,(y'-z)}\,\frac{(y'-z)\cdot (x-z)\,(x'-y')\cdot (y-w)}{(y'-z)^{2}\,(x'-y')^{2}\,(y-w)^{2}\,(x-z)^{2}}\,\mathcal C(y') &&(\delta_{jm}\delta_{ik'}),\notag\\
F_c(p^+)&=\int_{w,y',z,y,x,x'}\,e^{-ik\cdot (y'-w)}\, e^{-ik\cdot \frac{p^{+}}{k^{+}}\,(y'-z)}\,\frac{(y'-z)\cdot (y-w)\,(x'-y')\cdot (x-z)}{(y'-z)^{2}\,(x'-y')^{2}\,(y-w)^{2}\,(x-z)^{2}}\,\mathcal C(y') &&(\delta_{im}\delta_{jk'}).
\label{eq:F}
\end{align}
With $C=\frac{1}{(2\pi)^{3}}\, \frac{ig^{4}\,f^{c'd'a}}{4\pi^{4}}\, \frac{1}{k^{+}}$ and your $\left[\frac{1}{p^+}\right]_+f(p^+)=\frac{1}{p^+}\big(f(p^+)-f(0)\big)$, Eq.~\eqref{eq:userG3VI} is
\begin{align}
\text{line 1}&=C\int_{0}^{\vee-k^{+}}\frac{dp^{+}}{2\pi}\,\bigg[\frac{F_a(p^+)}{p^++k^+}-\frac{F_b(p^+)-F_b(0)}{p^+}-\frac{F_c(p^+)}{k^+}\bigg],\label{eq:line1}\\
\text{line 2}&=-\frac{C}{2\pi}\,F_b(0)\,\bigg[\log\Big(\frac{k^+}{\Lambda}\Big)+\log\Big(\frac{\vee}{k^+}-1\Big)\bigg],\qquad \int_\Lambda^{\vee-k^+}\frac{dp^+}{p^+}=\log\frac{k^+}{\Lambda}+\log\Big(\frac{\vee}{k^+}-1\Big).\label{eq:line2}
\end{align}
$F_b(0)$ has no $p^+$ phase: it is the integral of your line 2. ($\frac{C}{2\pi}=\frac{1}{(2\pi)^{3}}\frac{ig^{4}f^{c'd'a}}{8\pi^{5}}\frac{1}{k^{+}}$, the prefactor of your line 2.)

\section{Step 3: what the \texorpdfstring{$y'$}{y'} integral does at large \texorpdfstring{$p^+$}{p+}}
The two phases combine to one:
\begin{equation}
e^{-ik\cdot (y'-w)}\, e^{-ik\cdot \frac{p^{+}}{k^{+}}(y'-z)}=e^{-ik\cdot (z-w)}\;e^{-i\,t\,k\cdot(y'-z)},\qquad t=\frac{p^++k^+}{k^+},\quad 1\le t\le\frac{\vee}{k^+}.
\end{equation}
So the $y'$ integral is a Fourier transform at the momentum $q=t\,k$, which grows with $p^+$. All three integrands have the factor $\frac{(y'-z)^m}{(y'-z)^2}$, singular at $y'=z$; the rest, $G(y')$, is smooth there. The exact transform
\begin{equation}
\int d^2r\;e^{-iq\cdot r}\,\frac{r^m}{r^2}=-2\pi i\,\frac{q^m}{q^2}
\qquad\Big(\tfrac{r^m}{r^2}=\partial^m\tfrac12\log r^2,\ \int d^2r\,e^{-iq\cdot r}\log r^2=-\tfrac{4\pi}{q^2}\Big)
\label{eq:lemma0}
\end{equation}
decreases only like $\frac{1}{q}$. For large $q$ the oscillating phase cancels everything except $|y'-z|\lesssim\frac{1}{|q|}$, where $G(y')\simeq G(z)$:
\begin{equation}
\int_{y'}e^{-iq\cdot (y'-z)}\,\frac{(y'-z)^m}{(y'-z)^2}\,G(y')=-2\pi i\,\frac{q^m}{q^2}\,G(z)+O\Big(\frac{1}{q^2}\Big).
\label{eq:lemma}
\end{equation}
With $\frac{q^m}{q^2}=\frac1t\frac{k^m}{k^2}=\frac{k^+}{p^++k^+}\frac{k^m}{k^2}$, and $\int_{y'}e^{-ik\cdot(y'-w)}\frac{(y'-z)^m}{(y'-z)^2}=e^{-ik\cdot(z-w)}\big(-2\pi i\frac{k^m}{k^2}\big)$, this says
\begin{equation}
\boxed{F_n(p^+)=\frac{k^+}{p^++k^+}\,F_n^{z}+O\Big(\frac{k^{+2}}{p^{+2}}\Big)}\,,\qquad n=a,b,c,
\label{eq:Fz}
\end{equation}
where $F_n^z$ is $F_n$ at $p^+=0$ with $y'\to z$ everywhere except in $\frac{(y'-z)^m}{(y'-z)^2}$ and in $e^{-ik\cdot(y'-w)}$. For the third term:
\begin{equation}
F_c^{z}=\int_{w,y',z,y,x,x'}\, e^{-ik\cdot (y'-w)}\,\frac{(y'-z)\cdot (y-w)\,(x'-z)\cdot (x-z)}{(y'-z)^{2}\,(y-w)^{2}\,(x'-z)^{2}\,(x-z)^{2}}\;\mathcal C(z).
\label{eq:Fcz}
\end{equation}
\begin{itemize}
\item In words: the $w'$ integral set $w'=\frac{(p^++k^+)y'-p^+z}{k^+}$, i.e.\ $y'=\frac{k^+w'+p^+z}{p^++k^+}$. For $p^+\gg k^+$ the parent sits on $z$: the measured gluon ($k^+$) is emitted by the much faster gluon at $z$. This is the evolution of the projectile between $k^+$ and $\vee$, i.e.\ BFKL.
\item The second singular point, $y'\to x'$, also gives a $\frac{k^+}{p^++k^+}$ tail, but with the phase $e^{-i\frac{p^+}{k^+}k\cdot(x'-z)}$, which oscillates in $p^+$. Its $p^+$ integral converges, so it gives no logarithm.
\item Eq.~\eqref{eq:lemma} is exact for a Gaussian $G$ (Sec.~\ref{sec:toy}), and was checked numerically for a non-Gaussian $G$ in \texttt{BFKL\_final.pdf} (ratio $0.999996$, $1.000015$, $1.000054$, $1.000205$ at $q=6k,11k,21k,41k$).
\end{itemize}

\section{Step 4: which term of line 1 can grow like \texorpdfstring{$\log\frac{\vee}{k^+}$}{log(vee/k+)}}
Put Eq.~\eqref{eq:Fz} into Eq.~\eqref{eq:line1}. A $p^+$ integrand that falls like $\frac{1}{p^{+2}}$ gives a finite integral for $\vee\to\infty$. One that falls like $\frac{1}{p^+}$ gives $\log\frac{\vee}{k^+}$.
\begin{center}
\renewcommand{\arraystretch}{1.6}
\begin{tabular}{llll}
\toprule
term of line 1 & $p^+$ factor & integrand at $p^+\gg k^+$ & $\int^{\vee-k^+}dp^+$\\\midrule
$\delta_{k'm}\delta_{ij}$ & $\frac{1}{p^++k^+}$ & $\frac{k^+F_a^z}{(p^++k^+)^2}$ & finite\\
$\delta_{jm}\delta_{ik'}$, the $f(p^+)$ of $[\frac{1}{p^+}]_+$ & $-\frac{1}{p^+}$ & $-\frac{k^+F_b^z}{p^+(p^++k^+)}$ & finite\\
$\delta_{jm}\delta_{ik'}$, the $-f(0)$ of $[\frac{1}{p^+}]_+$ & $-\frac{1}{p^+}$ & $+\frac{F_b(0)}{p^+}$ & $+F_b(0)\log(\frac{\vee}{k^+}-1)+\dots$ \ (Step 5)\\
$\delta_{im}\delta_{jk'}$ & $-\frac{1}{k^+}$ & $-\frac{F_c^z}{p^++k^+}$ & $-F_c^z\log\frac{\vee}{k^+}+\dots$ \ (Step 6)\\\bottomrule
\end{tabular}
\end{center}
So two pieces of line 1 grow like $\log\frac{\vee}{k^+}$: the subtraction of the plus prescription, and the $-\frac{1}{k^+}$ term. They have different coefficients, $F_b(0)$ and $F_c^z$, and they do different things.

\section{Step 5: the \texorpdfstring{$\log(\frac{\vee}{k^+}-1)$}{log(vee/k+ - 1)} of line 2 is cancelled}
Split the plus-prescribed integral at $p^+=k^+$ (for $\vee>2k^+$) and use $\int_{k^+}^{\vee-k^+}\frac{dp^+}{p^+}=\log(\frac{\vee}{k^+}-1)$, exactly:
\begin{equation}
-\int_0^{\vee-k^+}dp^+\,\frac{F_b(p^+)-F_b(0)}{p^+}
=-\int_0^{k^+}dp^+\,\frac{F_b(p^+)-F_b(0)}{p^+}-\int_{k^+}^{\vee-k^+}dp^+\,\frac{F_b(p^+)}{p^+}+F_b(0)\log\Big(\frac{\vee}{k^+}-1\Big).
\end{equation}
Add line 2, Eq.~\eqref{eq:line2}: $\log(\frac{\vee}{k^+}-1)$ drops out,
\begin{equation}
\text{line 1}\big|_{\delta_{jm}\delta_{ik'}}+\text{line 2}=\frac{C}{2\pi}\bigg[-F_b(0)\log\frac{k^+}{\Lambda}-\int_0^{k^+}dp^+\,\frac{F_b(p^+)-F_b(0)}{p^+}-\int_{k^+}^{\vee-k^+}dp^+\,\frac{F_b(p^+)}{p^+}\bigg].
\label{eq:bcancel}
\end{equation}
The two integrals are finite for $\vee\to\infty$ ($F_b(p^+)\sim\frac{k^+}{p^+}$). What is left is JIMWLK ($\log\frac{k^+}{\Lambda}$) and a finite rest, with \textbf{no} $\log\frac{\vee}{k^+}$. Your $\log(\frac{\vee}{k^+}-1)$ only appears because $\left[\frac{1}{p^+}\right]_+$ subtracts $F_b(0)$ all the way up to $p^+=\vee-k^+$; the subtraction and line 2 cancel each other.

\section{Step 6: the hidden BFKL term (B) in the \texorpdfstring{$-\frac{1}{k^+}$}{-1/k+} term}
Since $\int_0^{\vee-k^+}\frac{dp^+}{p^++k^+}=\log\frac{\vee}{k^+}$ exactly,
\begin{equation}
-\frac{1}{k^+}\int_0^{\vee-k^+}dp^+\,F_c(p^+)=-F_c^z\,\log\frac{\vee}{k^+}\;-\;\frac{1}{k^+}\int_0^{\vee-k^+}dp^+\,\Big[F_c(p^+)-\frac{k^+}{p^++k^+}\,F_c^z\Big],
\label{eq:csplit}
\end{equation}
and by Eq.~\eqref{eq:Fz} the last integrand falls like $\frac{1}{p^{+2}}$, so the integral is finite for $\vee\to\infty$. The first term, times $\frac{C}{2\pi}$, is the BFKL term of Row VI. With Eq.~\eqref{eq:Fcz} (the contractions of $\delta_{im}\delta_{jk'}$: $(y'-z)\cdot(y-w)$ and $(x'-y')\cdot(x-z)\to(x'-z)\cdot(x-z)$; $U^{e'c'}(y')\to U^{e'c'}(z)$):
\begin{align}
\text{(B)}_{\rm VI}={}&-\frac{1}{(2\pi)^{3}}\, \frac{ig^{4}\,f^{c'd'a}}{8\pi^{5}}\, \frac{1}{k^{+}}\,\log{\left(\frac{\vee}{k^{+}}\right)}\int_{w,y',z,y,x,x'} e^{-ik\cdot (y'-w)}\,\frac{(y'-z)\cdot (y-w)\,(x'-z)\cdot (x-z)}{(y'-z)^{2}\,(y-w)^{2}\,(x'-z)^{2}\,(x-z)^{2}}\notag \\
& \left[U^{e'c'}(z)\,U^{d'b}(z)\,\rho^{e'}(x')-U^{bd'}(z)\,U^{c'e'}(x')\,\rho^{e'}(x')\right]\notag \\
& \times\left[U^{bc}(z)\,U^{ad}(y)\,U^{ce}(x)\,\rho^{d}(y)\,\rho^{e}(x)-U^{ad}(y)\,\rho^{d}(y)\,\rho^{b}(x)\right] \label{eq:BVI}
\end{align}
This is the term (B) of Group III Row VI in \texttt{BFKL\_rows.pdf}, Eq.~(4.8).

\section{Step 7: Row VI \texorpdfstring{$=$}{=} JIMWLK \texorpdfstring{$+$}{+} BFKL \texorpdfstring{$+$}{+} rest}
Collecting Steps 5 and 6 (the $\delta_{k'm}\delta_{ij}$ term is finite):
\begin{equation}
\text{Row VI}=\underbrace{-\frac{C}{2\pi}\,F_b(0)\,\log\frac{k^+}{\Lambda}}_{\text{JIMWLK (the }\log\frac{k^+}{\Lambda}\text{ of your line 2)}}\;+\;\underbrace{\text{(B)}_{\rm VI}}_{\text{BFKL}}\;+\;\text{Rest}_{\rm VI}.
\end{equation}
The rest can be written in two equivalent ways:
\begin{itemize}
\item \textbf{In your form} (as in \texttt{Rest\_terms.pdf}): line 1 as you have it, $+$ line 2 with $\log(\frac{\vee}{k^+}-1)$ only, $-$(B):
\begin{equation}
\text{Rest}_{\rm VI}=\text{line 1}-\frac{C}{2\pi}\,F_b(0)\log\Big(\frac{\vee}{k^+}-1\Big)+C\int_0^{\vee-k^+}\frac{dp^+}{2\pi}\,\frac{F_c^z}{p^++k^+}.
\label{eq:restyour}
\end{equation}
The last term is $-$(B), written with $\log\frac{\vee}{k^+}=\int_0^{\vee-k^+}\frac{dp^+}{p^++k^+}$. It is a subtraction at the hard end $p^+\to\vee$, as $[\frac{1}{p^+}]_+$ is one at the soft end $p^+\to0$.
\item \textbf{Manifestly finite} (Eqs.~\eqref{eq:bcancel}, \eqref{eq:csplit}): every integrand falls like $\frac{1}{p^{+2}}$, so one may even set $\vee\to\infty$:
\begin{equation}
\text{Rest}_{\rm VI}=\frac{C}{2\pi}\bigg\{\int_0^{\vee-k^+}dp^+\frac{F_a(p^+)}{p^++k^+}-\int_0^{k^+}dp^+\frac{F_b(p^+)-F_b(0)}{p^+}-\int_{k^+}^{\vee-k^+}dp^+\frac{F_b(p^+)}{p^+}-\frac{1}{k^+}\int_0^{\vee-k^+}dp^+\Big[F_c(p^+)-\frac{k^+F_c^z}{p^++k^+}\Big]\bigg\}.
\end{equation}
\end{itemize}

\section{Step 8: numbers (toy model)}\label{sec:toy}
To see Steps 4--6 with numbers, take $|k|=1$ and replace, in the $y'$ integral, everything except $\frac{(y'-z)^m}{(y'-z)^2}$ and the phase by $e^{-(y'-z)^2/2}$. Then Eq.~\eqref{eq:lemma} is exact:
\begin{equation}
\int d^2y'\,e^{-i\,t\,k\cdot(y'-z)}\,\frac{(y'-z)^m}{(y'-z)^2}\,e^{-(y'-z)^2/2}=-2\pi i\,\frac{k^m}{k^2}\;\frac{1-e^{-t^2/2}}{t},
\end{equation}
so $F(p^+)=F^z\,f(t)$ with $f(t)=\frac{1-e^{-t^2/2}}{t}\to\frac1t=\frac{k^+}{p^++k^+}$, $f(1)=F(0)/F^z=0.3935$. Use the same $F$ in all three terms (only the $p^+$ factor differs). In units of $\frac{C}{2\pi}F^z$ (\texttt{why\_toy.py}):
\begin{center}
\begin{tabular}{rrrrrrr}
\toprule
$\frac{\vee}{k^+}$ & $\delta_{k'm}\delta_{ij}$ & $[\frac{1}{p^+}]_+$ term & line 2: $\log(\frac{\vee}{k^+}-1)$ & sum of the two & $-\frac{1}{k^+}$ term & $-\frac{1}{k^+}$ term $+\log\frac{\vee}{k^+}$\\\midrule
$10^{1}$ & $+0.69116$ & $+0.17814$ & $-0.86454$ & $-0.68640$ & $-2.02270$ & $+0.27989$\\
$10^{2}$ & $+0.78116$ & $+1.02633$ & $-1.80804$ & $-0.78171$ & $-4.32528$ & $+0.27989$\\
$10^{3}$ & $+0.79016$ & $+1.92683$ & $-2.71760$ & $-0.79076$ & $-6.62787$ & $+0.27989$\\
$10^{4}$ & $+0.79106$ & $+2.83228$ & $-3.62395$ & $-0.79166$ & $-8.93045$ & $+0.27989$\\
$10^{6}$ & $+0.79116$ & $+4.64422$ & $-5.43598$ & $-0.79176$ & $-13.53562$ & $+0.27989$\\
\bottomrule
\end{tabular}
\end{center}
\begin{itemize}
\item The $\delta_{k'm}\delta_{ij}$ term converges.
\item The $[\frac{1}{p^+}]_+$ term grows by $f(1)\log10=0.9060$ per decade, and line 2 falls by the same amount; their sum converges (Step 5).
\item The $-\frac{1}{k^+}$ term falls by $\log10=2.3026$ per decade. It is $-\log\frac{\vee}{k^+}$ plus a constant, exactly: the coefficient of the logarithm is $F^z$, the $y'\to z$ value (Step 6). Adding $+\log\frac{\vee}{k^+}$, i.e.\ subtracting (B), leaves a constant.
\end{itemize}

\section{The other rows}
\paragraph{Group III Row III, first part.} The same three terms, with $\frac{ig^{4}f^{c'd'a}}{8\pi^{4}}$ and the colour factor with the anticommutators. Steps 2--7 are unchanged:
\begin{align}
\text{(B)}_{\rm III,1}={}&-\frac{1}{(2\pi)^{3}}\, \frac{ig^{4}\,f^{c'd'a}}{16\pi^{5}}\, \frac{1}{k^{+}}\,\log{\left(\frac{\vee}{k^{+}}\right)}\int_{w,x',y',z,y,x} e^{-ik\cdot (y'-w)}\,\frac{(y'-z)\cdot (y-w)\,(x'-z)\cdot (x-z)}{(y'-z)^{2}\,(y-w)^{2}\,(x'-z)^{2}\,(x-z)^{2}}\notag \\
& \left[U^{e'c'}(z)\,U^{d'b}(z)\,\rho^{e'}(x')-U^{bd'}(z)\,U^{c'e'}(x')\,\rho^{e'}(x')\right]\notag \\
& \times\left[U^{ac}(w)\, \{\rho^{c}(y),\rho^{b}(x)\}-U^{bc}(z)\,U^{ad}(y)\,U^{ce}(x)\,\{\rho^{d}(y),\rho^{e}(x)\}\right] \label{eq:BIII,1}
\end{align}
\paragraph{Group III Row V.} The same, with $\frac{ig^{4}f^{c'ca}}{4\pi^{4}}$ and its colour factor:
\begin{align}
\text{(B)}_{\rm V}={}&-\frac{1}{(2\pi)^{3}}\, \frac{ig^{4}\,f^{c'ca}}{8\pi^{5}}\, \frac{1}{k^{+}}\,\log{\left(\frac{\vee}{k^{+}}\right)}\int_{w,y',z,y,x,x'} e^{-ik\cdot (y'-w)}\,\frac{(y'-z)\cdot (y-w)\,(x'-z)\cdot (x-z)}{(y'-z)^{2}\,(y-w)^{2}\,(x'-z)^{2}\,(x-z)^{2}}\notag \\
& \left[U^{e'c'}(z)\,\rho^{e'}(x')-U^{c'e'}(x')\,\rho^{e'}(x')\right]\left[U^{ce}(x)\,U^{ad}(y)\,\rho^{e}(x)\rho^{d}(y)-U^{ad}(w)\,U^{ce}(x)\,\rho^{e}(x)\rho^{d}(y)\right] \label{eq:BV}
\end{align}
\paragraph{Group III Row III, second part.} Here the bracket has six terms and an extra factor $\frac{1}{p^{+}(x-z)^{2}+k^{+}(x-w)^{2}}\simeq\frac{1}{p^+(x-z)^2}$ at large $p^+$. With $F(p^+)\simeq\frac{k^+}{p^++k^+}F^z$:
\begin{center}
\renewcommand{\arraystretch}{1.5}
\begin{tabular}{lll}
\toprule
term of the bracket & integrand at $p^+\gg k^+$ & $\int^{\vee-k^+}dp^+$\\\midrule
$\frac{\delta_{k'm}d_\perp p^+}{2(z-w)^2(p^++k^+)^2}[\dots]$ & $\propto\frac{1}{p^{+3}}$ & finite\\
$\frac{p^+\delta_{k'm}}{k^+(p^++k^+)}[\dots]$ & $\propto\frac{1}{p^{+2}}$ & finite\\
$\frac{\delta_{k'm}}{p^++k^+}[\dots]$ & $\propto\frac{1}{p^{+3}}$ & finite\\
$-[\frac{1}{p^+}]_+[\dots]$, $f(p^+)$ & $\propto\frac{1}{p^{+3}}$ & finite\\
$-[\frac{1}{p^+}]_+[\dots]$, $-f(0)$ & $+\frac{f(0)}{p^+}$ & cancels the $\log(\frac{\vee}{k^+}-1)$ of line 2 (Step 5)\\
$-\frac{p^+}{k^{+2}}[\dots]$ & $-\frac{1}{k^+(x-z)^2}\,\frac{1}{p^++k^+}\,[\dots]F^z$ & $\log\frac{\vee}{k^+}$: (B) (Step 6)\\
$-\frac{1}{k^+}[\dots]$ & $\propto\frac{1}{p^{+2}}$ & finite\\\bottomrule
\end{tabular}
\end{center}
The $-\frac{p^+}{k^{+2}}$ term, with $\frac{p^+}{p^+(x-z)^2+k^+(x-w)^2}\to\frac{1}{(x-z)^2}$ and $(x'-y')\to(x'-z)$, $U^{e'c'}(y')\to U^{e'c'}(z)$, gives
\begin{align}
\text{(B)}_{\rm III,2}={}&-\frac{1}{(2\pi)^{3}}\,\frac{g^{4}}{8\pi^{5}}\,\frac{1}{k^{+}}\,\log{\left(\frac{\vee}{k^{+}}\right)}\int_{w,x',y',z,x} e^{-ik\cdot (y'-w)}\,\frac{(x'-z)\cdot (x-z)}{(x'-z)^{2}\,(x-z)^{2}}\notag \\
& \times\bigg[\frac{(y'-z)\cdot (z-w)}{(y'-z)^{2}\,(z-w)^{2}}-\frac{(y'-z)\cdot (x-w)}{2\,(y'-z)^{2}\,(x-w)^{2}}\bigg]\notag \\
& \left[f^{c'd'a}\,U^{e'c'}(z)\,U^{d'b}(z)\,\rho^{e'}(x')-f^{c'd'a}\, U^{bd'}(z)\,U^{c'e'}(x')\,\rho^{e'}(x')\right]\notag \\
& \times\left[f^{cbd}\,U^{ac}(w)\,\rho^{d}(x) -f^{acd}\,U^{bc}(z)\,U^{de}(x)\,\rho^{e}(x)\right] \label{eq:BIII,2}
\end{align}
\paragraph{Group III Row IV.} Only four $\rho$'s; nothing of this kind.
\paragraph{Group III Rows I, II.} Their $\log\frac{\vee}{k^+}$ is not hidden: it is inside the logarithms, e.g.\ $\log\frac{(\vee-k^{+})(x-z)^{2}+k^{+}(x-w)^{2}}{k^{+}(x-w)^{2}+\Lambda(x-z)^{2}}=\log\frac{\vee}{k^+}+\log\frac{(1-\frac{k^+}{\vee})(x-z)^2+\frac{k^+}{\vee}(x-w)^2}{(x-w)^2}+O(\Lambda)$, and $\log(\frac{\vee}{k^+}-1)$ there has no plus prescription to cancel it. It is taken out in \texttt{BFKL\_rows.pdf} (blue).

\paragraph{Group II Row IV (the same in your variable $\xi$).} With $\xi=\frac{k^+}{p^++k^+}$, $\frac{d\xi}{\xi(1-\xi)}=-\frac{dp^+}{p^+}$, your $\frac{\xi}{(1-\xi)_+}$ is your $[\frac{1}{p^+}]_+$, and its subtraction cancels the $\log(\frac{\vee}{k^+}-1)$ of line 2, as in Step 5. The term that grows is $\frac{1-\xi}{\xi}$: with $\frac{d\xi}{\xi^2}\frac{1-\xi}{\xi}$ and the phase $e^{-ik\cdot(y'-y)/\xi}$, the $-\frac{d_\perp}{2}\log\frac{(y'-y)^2\mu^2}{4}$ of your pole bracket has the transform $\int d^2\Delta\,e^{-iq\cdot\Delta}\log\Delta^2=-\frac{4\pi}{q^2}=-\frac{4\pi\xi^2}{k^2}$. The integrand is then $\propto\frac{d\xi}{\xi}$, from $y'\to y$ (and $y,y'\to z$), which gives $\int_{k^+/\vee}^1\frac{d\xi}{\xi}=\log\frac{\vee}{k^+}$ times your $H_4$ (line 2 of Group II Row IV in your \texttt{BFKL.tex}, which is correct). The constant parts of the bracket (pole, $\gamma_E$, $-1$) multiply functions that are smooth at $y'=y$; their transforms fall fast in $\frac1\xi$, so they give no logarithm. In \texttt{Rest\_terms.pdf} the row therefore has $-\int_{k^+/\vee}^1\frac{d\xi}{\xi}\,H_4$.

\section{Why (B) is needed: the checks}\label{sec:checks}
\begin{itemize}
\item \textbf{BFKL$\otimes$LO.} With the terms (B), the BFKL list equals the $\log\frac{\vee}{k^+}$ list of \texttt{LL\_terms.pdf} ($4.4\times10^{-17}$), and its two-$\rho$ terms are BFKL$\otimes$LO (0 of 109 colour structures differ, \texttt{../final\_check/render.py}). In the kernel-by-kernel comparison of \texttt{BFKL\_final.pdf} all 67 structures agree with (B); without (B), 9 of the 67 differ. These are exactly the kernels in which the measured gluon is emitted by $z$: $\mathcal K(x-w)\cdot\mathcal K(w'-z)$, $\mathcal K(x-w')\cdot\mathcal K(w-z)$ and $K(w,w',z)K(x,x',z)$.
\item \textbf{Three $\rho$'s.} The three-$\rho$ parts of the (B) terms cancel among themselves (Rows III and V: $\mp3.6368$, Row VI: $0$, in the table of \texttt{BFKL\_rows.pdf}, Sec.~6). The three-$\rho$ test therefore passes with or without them ($3.5\times10^{-16}$ / $3.6\times10^{-16}$); (B) matters for the two-$\rho$ terms.
\item \textbf{The rest is finite only with $-$(B).} Toy model, Sec.~\ref{sec:toy}: without $-$(B) the rest falls by $2.30\,F^z$ per decade of $\frac{\vee}{k^+}$.
\end{itemize}
\end{document}
